 \documentclass[a4paper,11pt]{book} %Formato de página
\usepackage[spanish,es-tabla]{babel} %Adaptación de LATEX al español (es-tabla hace que el pie de las tablas se denomine "tabla")
\usepackage{anysize} %Permite ajustar los margenes con el comando "\marginsize"
\usepackage{enumitem} %Paquete para personalizar los entronos de listas (enumerate, itemize and description)
\usepackage{graphicx} %Paquete para introducir imagenes (1 punto = 0.3528 mm)
\usepackage{subcaption} %Paquete para manejar subimagenes (listas de imagenes) 
\usepackage{caption} %Permite configurar las leyendas de figuras y tablas
\usepackage{multicol} % define un entorno multicols que compone texto en múltiples columnas (hasta un máximo de 10) 
\usepackage{float} %Permite controlar la posición exacta de los objetos flotantes (imagenes, tablas,etc...)
\usepackage{ulem} %Agrega multiples estilos para el el subrayado.
\usepackage{amsmath} %Añade comandos al entorno matemático de latex
\usepackage{mathtools}%Añade comandos al entorno matemático de latex
\usepackage{units} %Permite modificar las unidades e introducir fracciones tipo m/s de forma bonita.
\usepackage{amsfonts} %Agrega fuentes matemáticas y simbolos de utilidad
\usepackage{amsthm} %Permite escribir lemas, colorarios, teoremas, etc... de forma profesional
\usepackage{amssymb}
\usepackage{wrapfig} %Permite escribir texto alrededor de las imagenes
\usepackage{multirow} %Agrega elementos y aumenta la eficacia del entorno de tablas
\usepackage{fancyhdr} %Entorno para escribir pies de páginas 
\usepackage{setspace} %Permite cambiar el espacio entre lineas
\usepackage{appendix} %Permite agragar y modificar apendices
\usepackage{longtable} %Permite escribir tablas que ocupan más de una página
\usepackage{mdframed}  %Permite enmarcar y colorear
\usepackage{tocloft} % Paquete para personalizar el índice
\usepackage[table,xcdraw]{xcolor} %Paquete para colorear las tablas con la web online de tablas
\usepackage[most]{tcolorbox} %Proporciona un entorno para cuadros de texto enmarcados y coloreados con una línea de encabezado$
\usepackage[hidelinks]{hyperref} %Permite crear hipervinculos
\usepackage[hang,flushmargin]{footmisc} %proporciona medios para cambiar el diseño de las notas al pie
\usepackage{nicefrac, xfrac}
\usepackage{dsfont} % Para \mathds
\usepackage{bbm}    % Para \mathbbm
\usepackage{natbib} %Bibliografia estilo APS
\usepackage{braket} % Para escribir bra-ket de forma sencilla

%------------------------------------------------------------------------------------------------------------------------------------------
                     %WEB CON TODAS LAS BIBLIOTECAS Y TODA SU DOCUMENTACIÓN "https://ctan.org/"%
%------------------------------------------------------------------------------------------------------------------------------------------

%Personalizamos los hipervinculos
\hypersetup{colorlinks = true, linkcolor=black}
\hypersetup{colorlinks = true, urlcolor=blue}

% Personalización del índice
\setcounter{tocdepth}{1}
\renewcommand{\cftsecfont}{\color{black}} % Sección en negro
\renewcommand{\cftsubsecfont}{\color{black}} % Subsección en negro
\setlength{\cftbeforesecskip}{1.5em} % Espaciado antes de cada sección en el índice
\setlength{\cftbeforesubsecskip}{1.2em} % Espaciado antes de cada subsección

%Preparamos la forma de nuestras páginas y elegimos un estilo para los pies y cabezas de página
\marginsize{3cm}{3cm}{2.5cm}{2cm}
\pagestyle{fancy}
\setlength{\headheight}{15pt}
\fancyhf{} % limpia todo
\fancyhead[RE]{\leftmark}         % capítulo en páginas pares
\fancyhead[LO]{\rightmark}        % sección en páginas impares
\fancyfoot[C]{\thepage}           % número de página abajo centrado
%------------------------------------------------------------------------------------------------------------------------------------------

\begin{document}
\begin{titlepage}
	\centering
	\vspace*{\fill}
	{\fontsize{30pt}{36pt}\selectfont\bfseries Teoría Cuántica de Campos\par}
	\vspace{.5cm}
	{\scshape\LARGE Universidad Complutense de Madrid\par}
	\vspace{100pt}
    \resizebox{0.8\textwidth}{!}{$ \left(i\gamma^\mu \partial_\mu - m\right)\psi = 0$}
	\par
	\vspace{100pt}
	{\scshape\LARGE Sergio Figaredo Arce\par}
	\vspace{0.5cm}
	{\Large Curso 2025/2026\par}
	\vspace*{\fill}
\end{titlepage}
\tableofcontents % Genera el índice automáticamente
\newpage

%------------------------------------------------------------------------------------------------------------------------------------------
\chapter{Introducción}
\begin{quote}
\textit{``No existen sistemas reales de una sola partícula en la naturaleza, ni siquiera 
sistemas de pocas partículas. La existencia de pares virtuales y de fluctuaciones muestra que los días de los números fijos de partículas han terminado.''}
\begin{flushright}
\textit{Viki Weisskopf}
\end{flushright}
\end{quote}
\section{Descripción general}
\noindent
En la física clásica, la razón principal para introducir el concepto de campo es 
construir leyes de la Naturaleza que sean \textit{locales}. Las antiguas leyes de 
Coulomb y Newton implican ``acción a distancia''. Esto significa que la fuerza 
sentida por un electrón (o un planeta) cambia inmediatamente si un protón (o una 
estrella) distante se mueve. Esta situación es filosóficamente insatisfactoria. Más 
importante aún, también es experimentalmente incorrecta. Las teorías de campo de 
Maxwell y Einstein remedian la situación, con todas las interacciones mediadas de 
forma local por el campo.
\medskip\newline
El requisito de localidad sigue siendo una fuerte motivación para estudiar las 
teorías de campos en el mundo cuántico. Sin embargo, hay razones adicionales para 
tratar el campo cuántico como fundamental.
\medskip\newline
\textbf{Las partículas no son objetos indestructibles}, son, de hecho, en su 
mayoría efímeras y fugaces. Este hecho verificado experimentalmente fue predicho 
primero por Dirac, quien comprendió cómo la relatividad implica la necesidad de las 
antipartículas.
\medskip\newline
La presencia de multitud de partículas y antipartículas a distancias cortas nos dice 
que cualquier intento de escribir una versión relativista de la ecuación de 
Schrödinger de una partícula está condenado al fracaso. No existe ningún mecanismo 
en la mecánica cuántica no relativista estándar para tratar los cambios en el número 
de partículas. De hecho, cualquier intento ingenuo de construir una versión 
relativista de la ecuación de Schrödinger de una partícula se encuentra con serios 
problemas. (Las probabilidades negativas, las torres infinitas de estados de energía 
negativa o la ruptura de la causalidad son los problemas habituales que surgen.) 
\medskip\newline
En cada caso, este fracaso nos está diciendo que una vez que entramos en el régimen 
relativista, necesitamos un nuevo formalismo para tratar estados con un número no 
especificado de partículas. Este formalismo es la teoría cuántica de campos (TCC).
\newpage\noindent
\textbf{Todas las partículas del mismo tipo son iguales}. Esto suena bastante trivial. ¡Pero no lo es! Lo que quiero decir es que dos 
electrones son idénticos en todos los aspectos. Lo mismo ocurre con cualquier otra partícula fundamental.  Una explicación que podría ofrecerse es 
que hay un mar de ``sustancia'' de protón llenando el universo y cuando hacemos un 
protón de alguna manera metemos la mano en esta sustancia y de ella moldeamos un 
protón. Entonces no es sorprendente que los protones producidos en diferentes partes 
del universo sean idénticos: están hechos de la misma sustancia. Resulta que esto es 
aproximadamente lo que ocurre. La ``sustancia'' es el campo del protón o, si se mira 
con suficiente detalle, los campos del quark y del gluon.
\medskip\newline
En mecánica 
cuántica hay que introducir estas estadísticas a mano y,
para estar de acuerdo con el experimento, se deben elegir las estadísticas de Bose (sin signo negativo) para partículas de espín entero, y las estadísticas de Fermi (con signo negativo) para partículas de espín semientero. En la teoría cuántica de 
campos, esta relación entre espín y estadística no es algo que haya que introducir a 
mano. Es, más bien, una consecuencia del formalismo.
\subsection*{¿Qué es la Teoría Cuántica de Campos?}
\noindent
Habiendo explicado por qué la TCC es necesaria, deberíamos dar una definición de lo que realmente es. 
La pista está en el nombre: es la cuantización de un campo clásico.
\medskip\newline
En la mecánica cuántica estándar, se nos enseña a tomar los grados de libertad clásicos y promoverlos a operadores que 
actúan sobre un espacio de Hilbert. Las reglas para cuantizar un campo no son diferentes. Así, los grados de libertad básicos en la teoría cuántica de campos son \textit{funciones con valor de operador del espacio y el tiempo}. Esto significa que estamos tratando con un número infinito de grados de libertad, al menos uno por cada punto del espacio.
\medskip\newline
Resultará que las posibles interacciones en la teoría cuántica de campos están gobernadas por unos pocos principios básicos: localidad, simetría y el flujo del grupo de renormalización (el desacoplamiento de los fenómenos de corta distancia de la física a escalas mayores). Estas ideas hacen de la TCC un marco muy robusto.

\subsection*{¿Para qué sirve la Teoría Cuántica de Campos?}
\noindent
Como hemos destacado anteriormente, para cualquier sistema relativista es una necesidad. Pero también es una herramienta muy útil en 
sistemas no relativistas con muchas partículas. La teoría cuántica de campos ha tenido un gran impacto en la materia condensada, la física de altas energías, la cosmología, la gravedad cuántica y las matemáticas puras. Es literalmente el lenguaje en el que están escritas las leyes de la Naturaleza.
\newpage
\section{Unidades y Notación}
\noindent
La naturaleza nos presenta tres constantes fundamentales con dimensiones: la 
velocidad de la luz $c$, la constante de Planck (dividida por $2\pi$) $\hbar$, 
y la constante de Newton $G$. Sus dimensiones son:

\begin{align}
    [c] &= LT^{-1} \\
    [\hbar] &= L^2 M T^{-1} \\
    [G] &= L^3 M^{-1} T^{-2}
\end{align}
Usaremos unidades naturales $\hbar = c = 1$, donde las siguientes magnitudes 
tienen las mismas dimensiones:
\begin{equation}
[\text{longitud}] = [\text{tiempo}] = [\text{energía}]^{-1} = [\text{masa}]^{-1}
\end{equation}
Una relación muy útil es:
\begin{equation}
    \hbar c = 197.3269631(49) \ \text{MeV fm} \implies 
    25 \ \text{GeV}^{-2} \simeq 10^{-30} \ \text{m}^2 = 10 \ \text{mbarn}
    \tag{1.1}
\end{equation}
donde $1 \ \text{fm} = 10^{-15} \ \text{m}$ 
\subsection{Unidades Heaviside-Lorentz}
\noindent
En el Sistema Internacional (SI), las ecuaciones de Maxwell se escriben como
\begin{equation}
\left\{
\begin{aligned}
    \vec{\nabla}\cdot \vec{E} &= \frac{\rho}{\varepsilon_0}, \\
    \vec{\nabla}\cdot \vec{B} &= 0, \\
    \vec{\nabla}\times \vec{B} 
    &= \mu_0 \vec{J}
    + \frac{1}{c^2}\frac{\partial \vec{E}}{\partial t}, \\
    \vec{\nabla}\times \vec{E}
    &= -\frac{\partial \vec{B}}{\partial t}.
\end{aligned}
\right.
\end{equation}
donde
\begin{equation}
    c^2 = \frac{1}{\varepsilon_0 \mu_0},
    \qquad
    \alpha = \frac{e_{\text{SI}}^2}{4\pi \varepsilon_0 \hbar c}.
\end{equation}
La definición de las unidades gaussianas consiste en absorber el factor
$\varepsilon_0$ dentro de la definición de la carga.
\begin{equation}
    Q_G^2 = \frac{Q_{\text{SI}}^2}{4\pi\varepsilon_0}.
\end{equation}
Por tanto, en unidades gaussianas la constante de estructura fina queda definida por:
\begin{equation}
    \alpha = \frac{e_G^2}{\hbar c}.
\end{equation}
Las ecuaciones de Maxwell en unidades gaussianas toman la forma
\begin{equation}
\left\{
\begin{aligned}
    \vec{\nabla}\cdot \vec{E}_G &= 4\pi \rho_G, \\
    \vec{\nabla}\cdot \vec{B}_G &= 0, \\
    \vec{\nabla}\times \vec{B}_G 
    - \frac{1}{c}\frac{\partial \vec{E}_G}{\partial t}
    &= \frac{4\pi}{c}\vec{J}_G, \\
    \vec{\nabla}\times \vec{E}_G
    + \frac{1}{c}\frac{\partial \vec{B}_G}{\partial t}
    &= 0.
\end{aligned}
\right.
\end{equation}
Las unidades de Heaviside--Lorentz son una versión racionalizada de las
unidades gaussianas, es decir, una elección de unidades en la que se absorben
los factores $4\pi$ que aparecen en las ecuaciones de Maxwell gaussianas.
La relación entre los campos y las fuentes gaussianas y los campos y fuentes
de Heaviside--Lorentz es
\begin{equation}
    \vec{E}_{HL} = \frac{\vec{E}_G}{\sqrt{4\pi}},
    \qquad
    \vec{B}_{HL} = \frac{\vec{B}_G}{\sqrt{4\pi}},
\end{equation}
\begin{equation}
    \rho_{HL} = \sqrt{4\pi}\rho_G,
    \qquad
    \vec{J}_{HL} = \sqrt{4\pi}\vec{J}_G.
\end{equation}
Con estas definiciones, las ecuaciones de Maxwell quedan racionalizadas:
\begin{equation}
\left\{
\begin{aligned}
    \vec{\nabla}\cdot \vec{E}_{HL} &= \rho_{HL}, \\
    \vec{\nabla}\cdot \vec{B}_{HL} &= 0, \\
    \vec{\nabla}\times \vec{E}_{HL}
    + \frac{1}{c}\frac{\partial \vec{B}_{HL}}{\partial t}
    &= 0, \\
    \vec{\nabla}\times \vec{B}_{HL}
    - \frac{1}{c}\frac{\partial \vec{E}_{HL}}{\partial t}
    &= \frac{1}{c}\vec{J}_{HL}.
\end{aligned}
\right.
\end{equation}
Estas serán las unidades que utilizaremos, donde la constante de estructura fina y la unidad de carga quedan definidas por:
\begin{equation}
    \alpha
    =
    \frac{e^2}{4\pi \hbar c}
    =
    \frac{1}{137.03599911(46)}.
\end{equation}
\begin{equation}
    e = \sqrt{4\pi \alpha},
\end{equation}
\noindent
Todo esto nos permite definir el potencial de Coulomb entre dos cargas de la siguiente forma:
\begin{equation}
    Q_1 = e q_1,
    \qquad
    Q_2 = e q_2,
\end{equation}
es
\begin{equation}
    V(r)
    =
    \frac{Q_1 Q_2}{4\pi r}
    =
    \frac{q_1 q_2 \alpha}{r}.
\end{equation}
\section{Geometría Diferencial}
\noindent
En esta sección se presenta el marco geométrico y las convenciones necesarias para trabajar con objetos definidos sobre variedades. De esta forma, para definir en el espacio-tiempo de Minkowski, utilizaremos la siguiente convención de signos
\begin{equation}
    \eta_{\mu\nu} = g_{\mu\nu} = g^{\mu\nu}
    =
    \begin{pmatrix}
        1 & 0 & 0 & 0 \\
        0 & -1 & 0 & 0 \\
        0 & 0 & -1 & 0 \\
        0 & 0 & 0 & -1
    \end{pmatrix}.
\end{equation}\noindent
Por tanto, la métrica\footnote{La métrica permite definir productos escalares relativistas, subir y bajar
índices, y distinguir entre componentes contravariantes y covariantes. En
particular, satisface $g^{\mu\alpha}g_{\alpha\nu} = \delta^\mu_{\nu}$.} tiene signatura
\begin{equation}
    \eta_{\mu\nu} = (+,-,-,-).
\end{equation}
\newpage\noindent
Usaremos la convención de Einstein de suma sobre índices repetidos. Es decir,
si un índice aparece una vez arriba y una vez abajo, se entiende que se suma
sobre todos sus valores:
\begin{equation}
A_\mu B^\mu
=
\sum_{\mu=0}^{3} A_\mu B^\mu
=
g_{\mu\nu} A^\mu B^\nu
=
A^0B^0 - A^1B^1 - A^2B^2 - A^3B^3.
\end{equation}
\noindent
Los índices griegos toman valores espacio-temporales
\[
    \mu,\nu,\rho,\sigma,... = 0,1,2,3.
\]
mientras que los índices latinos denotan únicamente componentes espaciales
\[
    i,j,k, ... = 1,2,3.
\]
\subsection{Vectores contravariantes y covariantes}
\noindent
Un vector geométrico puede escribirse como combinación lineal de los vectores de una base. En coordenadas \(x^\mu\), la base natural del espacio tangente se denota por:
\begin{equation}
    \partial_\mu \equiv \frac{\partial}{\partial x^\mu}.
\end{equation}
Por tanto, un vector se escribe como
\begin{equation}
    V = V^\mu \partial_\mu.
\end{equation}
Aquí es importante distinguir entre el vector completo, sus componentes y los
vectores de la base. Las cantidades $V^\mu$ son las componentes contravariantes del vector, y por eso llevan índices superiores. En cambio, $\partial_\mu$ son los vectores de la base coordenada, y llevan índice inferior.
\medskip\newline
Bajo un cambio de coordenadas \(x^\mu \mapsto x'^\mu\), las componentes
contravariantes del vector se transforman como:
\begin{equation}
    V'^\mu
    =
    \frac{\partial x'^\mu}{\partial x^\nu} V^\nu.
\end{equation}
La base, sin embargo, se transforma de forma opuesta:
\begin{equation}
    \partial'_\mu
    =
    \frac{\partial x^\nu}{\partial x'^\mu}\partial_\nu.
\end{equation}
Por eso el objeto completo $V = V^\mu \partial_\mu$
es independiente del sistema de coordenadas elegido.
\medskip\newline
De forma dual, un covector se escribe como: $$\omega = \omega_\mu dx^\mu$$
donde las $\omega_\mu$ representan las componentes covariantes del covector, y $dx^\mu$ son los elementos de la base dual.
\medskip\newline
Siguiendo un razonamiento similar, las componentes covariantes y la base dual se transforman bajo \(x^\mu \mapsto x'^\mu\) siguiendo la siguiente regla:
\begin{equation}
    \omega'_\mu
    =
    \frac{\partial x^\nu}{\partial x'^\mu}\omega_\nu \qquad 
    dx'^\mu
    =
    \frac{\partial x'^\mu}{\partial x^\nu}dx^\nu.
\end{equation}
\newpage\noindent
La métrica permite pasar de componentes contravariantes a covariantes, y
viceversa:
\begin{equation}
    V_\mu = g_{\mu\nu}V^\nu,
    \qquad
    V^\mu = g^{\mu\nu}V_\nu.
\end{equation}
\noindent
\textbf{Ejemplo 1} 
\medskip\newline\noindent
Si $V^\mu = (V^0,V^1,V^2,V^3)$,
entonces, con la métrica de Minkowski,
\begin{equation}
    V_\mu
    =
    g_{\mu\nu}V^\nu
    =
    (V^0,-V^1,-V^2,-V^3).
\end{equation}
\begin{flushright}
    \(\square\)
\end{flushright}
Así, para el cuadrivector posición se tiene:
\begin{equation}
    x^\mu = (x^0,\vec{x}) = (t,\vec{x}),
    \qquad
    x_\mu = (t,-\vec{x}).
\end{equation}
\subsection{Derivadas parciales y operador d'Alembertiano}
\noindent
La derivada parcial, que mide cómo cambian numéricamente las componentes de un objeto, se denota por:
\begin{equation}
    \partial_\mu = \frac{\partial}{\partial x^\mu}.
\end{equation}
Al subir el índice con la métrica obtenemos $\partial^\mu = g^{\mu\nu}\partial_\nu$, que con nuestra convención de signos:
\begin{equation}
    \partial_\mu = (\partial_0,\partial_i),
    \qquad
    \partial^\mu = (\partial_0,-\partial_i).
\end{equation}
\noindent donde $\partial_0 = \frac{\partial}{\partial t},
    \partial_i = \frac{\partial}{\partial x^i}$.
\medskip\newline\noindent
El operador d'Alembertiano generaliza para un espacio-tiempo de Minkowski el operador laplaciano clásico, quedando definido como el producto escalar del cuadrigradiente consigo mismo:
\begin{equation}
    \Box
    =
    \partial_\mu \partial^\mu
    =
    \partial_0^2 - \nabla^2,
\end{equation}
donde $\nabla^2
    =
    \partial_x^2 + \partial_y^2 + \partial_z^2$.
\medskip\newline
De esta forma, el operador cuadrimomento se puede representar como:
\begin{equation}
    p^\mu = i\partial^\mu.
\end{equation}
\noindent
donde las componentes espaciales y temporal vienen dadas por:,
\begin{equation}
    p^0 = i\frac{\partial}{\partial x^0}
        = i\frac{\partial}{\partial t}
\qquad
    p^i = i\partial^i = -i\partial_i
        = -i\frac{\partial}{\partial x^i}.
\end{equation}
\subsection*{Derivadas covariantes}
\noindent
En un espacio plano, como el espacio de Minkowski escrito en coordenadas
cartesianas, las derivadas covariantes coinciden con las derivadas parciales.
Sin embargo, en coordenadas generales o en espacios curvos, es necesario
introducir términos de conexión para que las derivadas transformen
correctamente como tensores.
\medskip\newline
Para un vector contravariante, la derivada covariante se define como
\begin{equation}
    \nabla_\mu V^\nu
    =
    \partial_\mu V^\nu
    +
    \Gamma^\nu_{\mu\rho}V^\rho.
\end{equation}
\noindent
Para un covector, la derivada covariante es
\begin{equation}
    \nabla_\mu \omega_\nu
    =
    \partial_\mu \omega_\nu
    -
    \Gamma^\rho_{\mu\nu}\omega_\rho.
\end{equation}
\noindent
La conexión se denota mediante los símbolos de Christoffel, estos coeficientes indican cómo cambia la base de vectores de un punto a otro. En una variedad con métrica \(g_{\mu\nu}\), y usando la conexión de Levi-Civita, los símbolos de Christoffel vienen dados explícitamente por:
\begin{equation}
    \Gamma^\rho_{\mu\nu}
    =
    \frac{1}{2}g^{\rho\sigma}
    \left(
        \partial_\mu g_{\nu\sigma}
        +
        \partial_\nu g_{\mu\sigma}
        -
        \partial_\sigma g_{\mu\nu}
    \right).
\end{equation}
\noindent
Esta conexión es compatible con la métrica $\nabla_\rho g_{\mu\nu} = 0$, y no tiene torsión $\Gamma^\rho_{\mu\nu} = \Gamma^\rho_{\nu\mu}$.
\medskip\newline
Dado un tensor de tipo \((r,s)\), $T^{\alpha_1\cdots\alpha_r}{}_{\beta_1\cdots\beta_s}$, su derivada covariante queda definida según:
\begin{equation}
\begin{aligned}
    \nabla_\mu
    T^{\alpha_1\cdots\alpha_r}{}_{\beta_1\cdots\beta_s}
    &=
    \partial_\mu
    T^{\alpha_1\cdots\alpha_r}{}_{\beta_1\cdots\beta_s}
    \\
    &\quad
    +
    \sum_{k=1}^{r}
    \Gamma^{\alpha_k}_{\mu\rho}
    T^{\alpha_1\cdots \rho \cdots \alpha_r}{}_{\beta_1\cdots\beta_s}
    \\
    &\quad
    -
    \sum_{k=1}^{s}
    \Gamma^\rho_{\mu\beta_k}
    T^{\alpha_1\cdots\alpha_r}{}_{\beta_1\cdots \rho \cdots \beta_s}.
\end{aligned}
\end{equation}
donde en la primera suma se sustituye el índice contravariante
\(\alpha_k\) por \(\rho\), mientras que en la segunda suma se sustituye el
índice covariante \(\beta_k\) por \(\rho\).
\subsection*{Derivada exterior}
\noindent
La derivada exterior es una operación que actúa sobre formas diferenciales.
A diferencia de la derivada covariante, no depende de la métrica ni de los
símbolos de Christoffel. 
\medskip\newline
Si \(f\) es una función escalar, es decir, una \(0\)-forma, su derivada exterior es:
\begin{equation}
    df = \partial_\mu f\, dx^\mu.
\end{equation}
\noindent
Si \(\omega\) es una \(1\)-forma, $\omega = \omega_\mu dx^\mu$, su derivada exterior es la \(2\)-forma
\begin{equation}
    d\omega
    =
    \partial_\mu \omega_\nu \, dx^\mu \wedge dx^\nu.
\end{equation}
\noindent
Debido a la antisimetría del producto exterior,
\begin{equation}
    dx^\mu \wedge dx^\nu = - dx^\nu \wedge dx^\mu,
\end{equation}
también podemos escribir
\begin{equation}
    d\omega
    =
    \frac{1}{2}
    \left(
        \partial_\mu \omega_\nu
        -
        \partial_\nu \omega_\mu
    \right)
    dx^\mu \wedge dx^\nu.
\end{equation}
La componente de \(d\omega\) es, por tanto,
\begin{equation}
    (d\omega)_{\mu\nu}
    =
    \partial_\mu \omega_\nu
    -
    \partial_\nu \omega_\mu.
\end{equation}

Si escribimos esta expresión usando derivadas covariantes, obtenemos
\begin{equation}
    (d\omega)_{\mu\nu}
    =
    \nabla_\mu \omega_\nu
    -
    \nabla_\nu \omega_\mu.
\end{equation}
Sustituyendo explícitamente los términos de conexión,
\begin{equation}
\begin{aligned}
    \nabla_\mu \omega_\nu
    -
    \nabla_\nu \omega_\mu
    &=
    \left(
        \partial_\mu \omega_\nu
        -
        \Gamma^\rho_{\mu\nu}\omega_\rho
    \right)
    -
    \left(
        \partial_\nu \omega_\mu
        -
        \Gamma^\rho_{\nu\mu}\omega_\rho
    \right)
    \\
    &=
    \partial_\mu \omega_\nu
    -
    \partial_\nu \omega_\mu
    -
    \Gamma^\rho_{\mu\nu}\omega_\rho
    \Gamma^\rho_{\nu\mu}\omega_\rho.
\end{aligned}
\end{equation}
\noindent
Para la conexión de Levi-Civita, que es simétrica en los índices inferiores, los términos de conexión se cancelan. Por tanto:
\begin{equation}
    (d\omega)_{\mu\nu}
    =
    \partial_\mu \omega_\nu
    -
    \partial_\nu \omega_\mu.
\end{equation}
Esto muestra que la derivada exterior de una \(1\)-forma no depende de la
conexión, sino únicamente de la parte antisimétrica de sus derivadas parciales.
\medskip\newline
Para una \(p\)-forma general, $ \omega
    =
    \frac{1}{p!}
    \omega_{\mu_1\cdots\mu_p}
    dx^{\mu_1}\wedge\cdots\wedge dx^{\mu_p}$,la derivada exterior se define como:
\begin{equation}
    d\omega
    =
    \frac{1}{p!}
    \partial_\nu \omega_{\mu_1\cdots\mu_p}
    dx^\nu \wedge dx^{\mu_1}\wedge\cdots\wedge dx^{\mu_p}.
\end{equation}
Equivalentemente, en componentes,
\begin{equation}
    (d\omega)_{\nu\mu_1\cdots\mu_p}
    =
    (p+1)\,
    \partial_{[\nu}
    \omega_{\mu_1\cdots\mu_p]},
\end{equation}
donde los corchetes indican antisimetrización\footnote{La antisimetrización de índices consiste en tomar una expresión tensorial y quedarse únicamente con la parte que cambia de signo al intercambiar dos índices. Para dos índices, como
\begin{equation}
    T_{[\mu\nu]}
    \equiv
    \frac{1}{2}
    \left(
  T_{\mu\nu}
   -
T_{\nu\mu}
    \right).
\end{equation}} de los índices.
\subsection*{Relación entre derivada covariante y derivada exterior}
\noindent
La derivada covariante mide cómo cambia un tensor teniendo en cuenta la
variación de la base. Por eso depende de la conexión
\medskip\newline
En cambio, la derivada exterior mide la variación antisimétrica de una forma
diferencial y no necesita conexión. La razón es que, al antisimetrizar, los términos simétricos de la conexión de Levi-Civita se cancelan.
\subsection{Producto escalar relativista}
\noindent
El producto escalar entre dos cuadrivectores se define mediante la métrica:
\begin{equation}
    V\cdot W
    =
    g_{\mu\nu}V^\mu W^\nu
    =
    V_\nu W^\nu.
\end{equation}
\noindent
En particular, para el cuadrimomento,
\begin{equation}
    p^\mu = i\partial^\mu = (p^0,\vec{p}) = (E,\vec{p}),
\qquad
    p_\mu = g_{\mu\nu}p^\nu = (E,-\vec{p}).
\end{equation}
\noindent
Lo que nos permite discernir el siguiente resultado:
\begin{equation}
    p^\mu p_\mu
    =
    p^\mu g_{\mu\nu}p^\nu
    =
    E^2 - |\vec{p}|^2.
\end{equation}
\noindent
Para una partícula de masa \(m\), esta cantidad es invariante y satisface
\begin{equation}
    p^\mu p_\mu = m^2.
\end{equation}
\section{Matrices de Dirac}
\noindent
Las matrices de Dirac \(\gamma^\mu\) aparecen al construir ecuaciones
relativistas de primer orden en derivadas, como la ecuación de Dirac. Estas
matrices satisfacen el álgebra de Clifford asociada a la métrica de Minkowski,
\begin{equation}
    \left\{
        \gamma^\mu,\gamma^\nu
    \right\}
    =
    \gamma^\mu\gamma^\nu
    +
    \gamma^\nu\gamma^\mu
    =
    2\eta^{\mu\nu}.
\end{equation}
\noindent
Con la convención de signatura se obtiene:
\begin{equation}
    \left(\gamma^0\right)^2 = 1,
    \qquad
    \left(\gamma^i\right)^2 = -1,
    \qquad i=1,2,3.
\end{equation}
\noindent
Además, \(\gamma^0\) es hermítica, mientras que las matrices espaciales
\(\gamma^i\) son antihermíticas:

\begin{equation}
    \left(\gamma^0\right)^\dagger = \gamma^0,
    \qquad
    \left(\gamma^i\right)^\dagger = -\gamma^i.
\end{equation}
\noindent
De forma compacta, esto puede escribirse como
\begin{equation}
    \left(\gamma^\mu\right)^\dagger
    =
    \gamma^0\gamma^\mu\gamma^0.
\end{equation}
\subsection{Matriz \texorpdfstring{\(\gamma^5\)}{gamma5}}
\noindent
La matriz \(\gamma^5\) se define como:
\begin{equation}
    \gamma^5
    =
    i\gamma^0\gamma^1\gamma^2\gamma^3.
\end{equation}
\noindent
Esta matriz satisface
\begin{equation}
    \left(\gamma^5\right)^2 = 1,
    \qquad
    \left(\gamma^5\right)^\dagger = \gamma^5,
    \qquad
    \left\{\gamma^5,\gamma^\mu\right\} = 0.
\end{equation}
\noindent
Es decir, \(\gamma^5\) es hermítica, su cuadrado es la identidad y anticonmuta
con todas las matrices \(\gamma^\mu\). Esta matriz será fundamental para
describir quiralidad y proyectores quirales.

\subsection{Generadores espinoriales de Lorentz}
\noindent
Se define también el tensor
\begin{equation}
    \sigma^{\mu\nu}
    =
    \frac{i}{2}
    \left[
        \gamma^\mu,\gamma^\nu
    \right],
\end{equation}
\noindent
donde

\begin{equation}
    \left[
        \gamma^\mu,\gamma^\nu
    \right]
    =
    \gamma^\mu\gamma^\nu
    -
    \gamma^\nu\gamma^\mu.
\end{equation}
\noindent
Los objetos \(\sigma^{\mu\nu}\) actuarán como generadores de
las transformaciones de Lorentz sobre espinores.
\subsection{Representación quiral o de Weyl}
\noindent
Una representación especialmente útil del álgebra de matrices \(\gamma^\mu\)
es la representación quiral, también llamada representación de Weyl. En ella:
\begin{equation}
    \gamma^0
    =
    \begin{pmatrix}
        0 & \mathbb{1} \\
        \mathbb{1} & 0
    \end{pmatrix},
    \qquad
    \gamma^i
    =
    \begin{pmatrix}
        0 & \sigma^i \\
        -\sigma^i & 0
    \end{pmatrix},
    \qquad
    \gamma^5
    =
    \begin{pmatrix}
        -\mathbb{1} & 0 \\
        0 & \mathbb{1}
    \end{pmatrix}.
\end{equation}
\noindent
Aquí \(\mathbb{1}\) representa la matriz identidad \(2\times 2\), y
\(\sigma^i\) son las matrices de Pauli. Esta representación separa de forma
natural las componentes quirales izquierda y derecha de un espinor.

\subsection{Representación estándar}
\noindent
Otra representación habitual es la representación ordinaria o estándar. En
esta base:
\begin{equation}
    \gamma^0
    =
    \begin{pmatrix}
        \mathbb{1} & 0 \\
        0 & -\mathbb{1}
    \end{pmatrix},
    \qquad
    \gamma^i
    =
    \begin{pmatrix}
        0 & \sigma^i \\
        -\sigma^i & 0
    \end{pmatrix},
    \qquad
    \gamma^5
    =
    \begin{pmatrix}
        0 & \mathbb{1} \\
        \mathbb{1} & 0
    \end{pmatrix}.
\end{equation}
Esta representación es útil en muchos cálculos explícitos porque separa de
forma clara las componentes de energía positiva y negativa en el límite no
relativista.
\subsection{Matrices sigma de Weyl}
\noindent
También se definen los objetos:

\begin{equation}
    \sigma^\mu
    =
    \left(
        \mathbb{1},
        \vec{\sigma}
    \right),
    \qquad
    \bar{\sigma}^\mu
    =
    \left(
        \mathbb{1},
        -\vec{\sigma}
    \right).
\end{equation}
\noindent
Estas matrices son útiles para escribir espinores de Weyl y para separar la
estructura quiral de las ecuaciones relativistas.

\subsection{Trazas útiles de matrices gamma}
\noindent
En el cálculo de secciones eficaces y tasas de desintegración aparecen con
frecuencia trazas de productos de matrices gamma. Algunas identidades básicas
son:
\begin{equation}
    \mathrm{Tr}
    \left(
        \gamma^\mu\gamma^\nu
    \right)
    =
    4\eta^{\mu\nu}.
\end{equation}
\noindent
Para cuatro matrices gamma:
\begin{equation}
    \mathrm{Tr}
    \left(
        \gamma^\mu\gamma^\nu\gamma^\rho\gamma^\sigma
    \right)
    =
    4
    \left(
        \eta^{\mu\nu}\eta^{\rho\sigma}
        -
        \eta^{\mu\rho}\eta^{\nu\sigma}
        +
        \eta^{\mu\sigma}\eta^{\nu\rho}
    \right).
\end{equation}
\noindent
Y con una matriz \(\gamma^5\):

\begin{equation}
    \mathrm{Tr}
    \left(
        \gamma^5
        \gamma^\mu
        \gamma^\nu
        \gamma^\rho
        \gamma^\sigma
    \right)
    =
    -4i\epsilon^{\mu\nu\rho\sigma}.
\end{equation}
\section{Matrices de Pauli}
\noindent
Las matrices de Pauli se definen como:
\begin{equation}
    \sigma^1
    =
    \begin{pmatrix}
        0 & 1 \\
        1 & 0
    \end{pmatrix},
    \qquad
    \sigma^2
    =
    \begin{pmatrix}
        0 & -i \\
        i & 0
    \end{pmatrix},
    \qquad
    \sigma^3
    =
    \begin{pmatrix}
        1 & 0 \\
        0 & -1
    \end{pmatrix}.
\end{equation}
\noindent
Estas matrices satisfacen la identidad
\begin{equation}
    \sigma^i\sigma^j
    =
    \delta^{ij}\mathbb{1}
    +
    i\epsilon^{ijk}\sigma^k.
\end{equation}
\noindent
Esta relación recoge tanto su parte simétrica:

\begin{equation}
    \left\{
        \sigma^i,\sigma^j
    \right\}
    =
    2\delta^{ij}\mathbb{1},
\end{equation}
\noindent:
como su parte antisimétrica:
\begin{equation}
    \left[
        \sigma^i,\sigma^j
    \right]
    =
    2i\epsilon^{ijk}\sigma^k.
\end{equation}
\newpage
\section{Problemas MC Relativista}
\noindent
En mecánica cuántica no relativista, el estado de una partícula
$\ket{\psi}$ evoluciona temporalmente mediante la ecuación de Schrödinger,
\begin{equation}
    i\hbar \frac{\partial}{\partial t}\psi(\vec{x},t)
    =
    H\psi(\vec{x},t),
\end{equation}
\noindent
donde
\begin{equation}
    \psi(\vec{x},t)
    =
    \braket{\vec{x}}{\psi(t)}
\end{equation}
\noindent
es la función de onda en la representación de posiciones y \(H\) es el
Hamiltoniano.
\medskip\newline
Para una teoría libre no relativista, la energía viene dada por
\begin{equation}
    E = \frac{\vec{p}^{\,2}}{2m},
    \qquad
    \vec{p} = -i\hbar \vec{\nabla}.
\end{equation}
\noindent
Sustituyendo el operador momento en la relación de energía se obtiene la
ecuación de Schrödinger libre:
\begin{equation}
    i\hbar \frac{\partial \psi}{\partial t}
    =
    -\frac{\hbar^2}{2m}\nabla^2\psi.
\end{equation}
En este caso, la densidad $\rho(\vec{x},t)=|\psi(\vec{x},t)|^2$ puede interpretarse como una densidad de probabilidad de encontrar la partícula
en el intervalo \([\vec{x},\vec{x}+d\vec{x}]\). 
\medskip\newline
Además, satisface una ecuación de continuidad:
\begin{equation}
    \frac{\partial \rho}{\partial t}
    =
    -\vec{\nabla}\cdot \vec{j},
\end{equation}
que expresa la conservación de la probabilidad, donde la corriente de probabilidad es:

\begin{equation}
    \vec{j}
    =
    \frac{i\hbar}{2m}
    \left(
        \psi \vec{\nabla}\psi^*
        -
        \psi^*\vec{\nabla}\psi
    \right).
\end{equation}
\subsection{Densidades de probabilidad negativas}
\noindent
El problema aparece cuando construimos una teoría relativista de una sola partícula. La relación relativista entre energía y momento es
\begin{equation}
    E^2 = \vec{p}^{\,2}c^2 + m^2c^4.
\end{equation}
\noindent
Promoviendo la energía y el momento a operadores,
\begin{equation}
    E \rightarrow i\hbar \frac{\partial}{\partial t},
    \qquad
    \vec{p} \rightarrow -i\hbar\vec{\nabla},
\end{equation}
\noindent
se obtiene
\begin{equation}
    -\hbar^2
    \frac{\partial^2 \psi}{\partial t^2}
    =
    \left(
        -\hbar^2 c^2\nabla^2
        +
        m^2c^4
    \right)\psi.
\end{equation}
\noindent
Equivalentemente,

\begin{equation}
    \frac{1}{c^2}
    \frac{\partial^2 \psi}{\partial t^2}
    =
    \nabla^2\psi
    -
    \frac{m^2c^2}{\hbar^2}\psi.
\end{equation}
\newpage\noindent
En unidades naturales, esta ecuación toma la forma:
\begin{equation}
    \left(
        \Box + m^2
    \right)\psi = 0,
\end{equation}
\noindent
que es la ecuación de Klein--Gordon.
\medskip\newline
Aunque esta ecuación es relativista, no permite interpretar directamente
\(\rho\) como una densidad de probabilidad positiva. La cantidad conservada que
se obtiene a partir de la ecuación de Klein--Gordon es:

\begin{equation}
    \rho
    =
    \frac{i\hbar}{2mc^2}
    \left(
        \psi^*\frac{\partial \psi}{\partial t}
        -
        \frac{\partial \psi^*}{\partial t}\psi
    \right).
\end{equation}
\noindent
Esta densidad satisface de nuevo una ecuación de continuidad,

\begin{equation}
    \frac{\partial \rho}{\partial t}
    =
    -\vec{\nabla}\cdot \vec{j}.
\end{equation}
\noindent
Sin embargo, a diferencia del caso no relativista, esta \(\rho\) puede tomar
cualquier signo. Es decir,

\begin{equation}
    \rho \in \mathbb{R},
    \qquad
    \rho \not\geq 0.
\end{equation}
\noindent
Por tanto, no puede interpretarse como una densidad de probabilidad. La razón profunda es que la ecuación de Klein--Gordon es una ecuación
diferencial de segundo orden en el tiempo. Por ello, las condiciones iniciales
\(\psi\) y \(\partial\psi/\partial t\) son independientes, lo que permite que
la densidad \(\rho\), aunque sea real, pueda ser negativa.

\subsection{Energías negativas}
\noindent
Otro problema característico de las ecuaciones de onda relativistas es la
aparición de soluciones con energías negativas. En el caso no relativista, las ecuaciones de onda libres solo tienen soluciones
de energía positiva. Por ejemplo, para la ecuación de Schrödinger,

\begin{equation}
    i\hbar \frac{\partial \psi}{\partial t}
    =
    -\frac{\hbar^2}{2m}\nabla^2\psi,
\end{equation}
\noindent
podemos tomar soluciones de onda plana de la forma
\begin{equation}
    \psi(\vec{x},t)
    \sim
    e^{-i(Et-\vec{p}\cdot\vec{x})/\hbar}
    =
    e^{-i(\omega t-\vec{k}\cdot\vec{x})},
\end{equation}
\noindent
con $
    E = \frac{\vec{p}^{\,2}}{2m}$. Como \(\vec{p}^{\,2}\geq 0\), se tiene

\begin{equation}
    E > 0.
\end{equation}
\noindent
En cambio, para una ecuación relativista libre, como la de Klein--Gordon en
unidades naturales, la relación de dispersión es:
\begin{equation}
    E^2 = \vec{p}^{\,2} + m^2.
\end{equation}
Por tanto:

\begin{equation}
    E = \sqrt{\vec{p}^{\,2}+m^2}.
\end{equation}
\newpage\noindent
Como las soluciones pueden aparecer con ambos signos de la energía:

\begin{equation}
\begin{cases}
    \psi_+
    \sim
    e^{-ip_\mu x^\mu}
    =
    e^{-i(Et-\vec{p}\cdot\vec{x})},
    \\[0.2cm]
    \psi_-
    \sim
    e^{+ip_\mu x^\mu}
    =
    e^{+i(Et-\vec{p}\cdot\vec{x})}.
\end{cases}
\end{equation}
\noindent
Actuando con el operador energía \(i\partial_t\), se obtiene

\begin{equation}
    i\frac{\partial \psi_+}{\partial t}
    =
    +E\psi_+,
\end{equation}
\noindent
mientras que
\begin{equation}
    i\frac{\partial \psi_-}{\partial t}
    =
    -E\psi_-.
\end{equation}
\noindent
Por tanto:
\begin{equation}
    \psi_+
    \quad \text{tiene energía} \quad +E>0,
\qquad
    \psi_-
    \quad \text{tiene energía} \quad -E<0.
\end{equation}
\noindent
Estos problemas indican que la interpretación correcta no debe hacerse en
términos de una función de onda relativista de una sola partícula, sino mediante teoría cuántica de campos, donde las soluciones de energía negativa se reinterpretan en términos de antipartículas y los campos, no las funciones de onda, son los objetos fundamentales.

\chapter{Simetrías de Lorentz y Poincarè}
\noindent
\textit{Mencionamos en la Introducción que la teoría cuántica de campos es una 
síntesis de los principios de la mecánica cuántica y de la relatividad especial. 
Nuestra primera tarea será entender cómo se implementa la simetría de Lorentz en 
la teoría de campos. 
\medskip\newline
Estudiaremos las representaciones del grupo de Lorentz en 
términos de campos e introduciremos los campos escalar, espinorial y vectorial. 
Examinaremos después la información proveniente de la invariancia de Poincaré. 
\medskip\newline
Este capítulo es bastante matemático y formal, pero una comprensión de este enfoque en teoría de grupos simplifica enormemente la construcción de los Lagrangianos para los distintos campos y proporciona en general una comprensión más profunda de diversos 
aspectos de la TCC.}


\section{Grupos y Grupos de Lie}
\noindent
Antes de estudiar simetrías continuas, como las transformaciones de Lorentz o Poincaré, conviene recordar qué es un grupo y cómo actúa sobre espacios
vectoriales. 
\subsection{Grupo}
\noindent
Un grupo es un conjunto de elementos \(G\), no necesariamente numerable,
dotado de una ley de composición interna $\cdot : G \times G \longrightarrow G$, tal que, para cualesquiera \(a,b,c \in G\), se satisfacen las siguientes propiedades:
\begin{enumerate}
    \item \textbf{Cierre:}
    \begin{equation}
        a,b \in G
        \quad \Longrightarrow \quad
        a\cdot b \in G.
    \end{equation}
    \item \textbf{Asociatividad:}
    \begin{equation}
        a\cdot (b\cdot c)
        =
        (a\cdot b)\cdot c.
    \end{equation}
    \item \textbf{Elemento neutro:} existe un elemento \(e \in G\) tal que
    \begin{equation}
        a\cdot e = e\cdot a = a,
        \qquad
        \forall a \in G.
    \end{equation}
    \item \textbf{Elemento inverso:} para cada \(a \in G\), existe un elemento
    \(a^{-1}\in G\) tal que
    \begin{equation}
        a\cdot a^{-1}
        =
        a^{-1}\cdot a
        =
        e.
    \end{equation}
\end{enumerate}
\noindent
La operación del grupo no tiene por qué ser conmutativa. Es decir, en general $a\cdot b \neq b\cdot a$. Cuando sí se cumple \(a\cdot b=b\cdot a\) para todo par de elementos, el grupo se llama \textit{abeliano}.
\newpage
\subsection{Grupos de Lie}
\noindent
Un grupo de Lie es un grupo cuyos elementos dependen de forma continua y
diferenciable de un conjunto de parámetros reales. Es decir, si $\theta = (\theta_1,\ldots,\theta_N)$ es un conjunto de \(N\) parámetros reales, podemos escribir sus elementos como $g(\theta) \in G$. 
\medskip\newline
El elemento neutro se obtiene para un valor particular de los parámetros, que normalmente se toma como \(\theta=0\):
\begin{equation}
    g(0)=e.
\end{equation}
\noindent
Además, el inverso de un elemento cercano a la identidad puede escribirse como
\begin{equation}
    g^{-1}(\theta)=g(-\theta).
\end{equation}
\noindent
La dimensión del grupo de Lie es el número de parámetros independientes que
lo describen. En este caso, $\dim G = N.$
\subsubsection{Subgrupo}
\noindent
Un subgrupo \(H\) de un grupo \(G\) es un subconjunto de \(G\) que también tiene estructura de grupo con la misma ley de composición. Es decir, si $H \subset G$, y \(h_1,h_2 \in H\), entonces:
\begin{equation}
    h_1 h_2 \in H,
    \qquad
    h^{-1}_1 \in H.
\end{equation}
\subsubsection{Subgrupo Invariante}
\noindent
Un subgrupo \(H\) se dice invariante, o normal, si para todo \(h\in H\) y todo \(g\in G\), se cumple:
\begin{equation}
    g h g^{-1} \in H.
\end{equation}
\noindent
Esto significa que, al conjugar un elemento de \(H\) por cualquier elemento del grupo \(G\), el resultado sigue perteneciendo a \(H\).
\subsubsection{Grupo Simple}
\noindent
Un grupo simple es aquel que no contiene ningún subgrupo invariante propio. Es decir, sus únicos subgrupos invariantes son el subgrupo trivial y el propio grupo:
\begin{equation}
    H = \{e\},
    \qquad
    \text{o}
    \qquad
    H = G.
\end{equation}
\noindent
Por tanto, un grupo simple no puede descomponerse en partes normales más pequeñas.
\section{Representaciones}
\noindent
Una representación de un grupo consiste en asociar a cada elemento del grupo
un operador lineal que actué sobre un espacio vectorial. Más precisamente, una representación \(D\) de un grupo \(G\) sobre un espacio vectorial \(V\) es una aplicación $D_R : G \longrightarrow GL(V)$, tal que:
\begin{equation}
    D_R(g_1 g_2)
    =
    D_R(g_1)D_R(g_2),
\end{equation}

y
\begin{equation}
    D_R(e)=\mathds{1}.
\end{equation}
\noindent
Aquí \(GL(V)\) es el grupo de transformaciones lineales invertibles sobre
\(V\), y suele denominarse base de la representación $R$. Si \(V\) tiene dimensión \(n\), entonces \(D_R(g)\) puede representarse como una matriz \(n\times n\).
\medskip\newline
La representación permite describir cómo actúan las simetrías del grupo sobre objetos matemáticos concretos, como vectores, campos, espinores o tensores. Una misma simetría abstracta puede actuar sobre objetos distintos de formas diferentes. Por eso es necesario distinguir entre el grupo abstracto y sus distintas representaciones.
\subsection{Representaciones equivalentes}
\noindent
Dos representaciones \(D_R\) y \(D_{R'}\) se dice que son equivalentes si están
relacionadas mediante un cambio de base. Es decir, si existe una matriz
invertible \(S\) tal que
\begin{equation}
    D_{R'}(g)
    =
    S^{-1}D_R(g)S,
    \qquad
    \forall g\in G.
\end{equation}
\noindent
En ese caso, ambas representaciones describen la misma acción del grupo, pero
escrita en bases distintas del espacio vectorial. Por tanto, no representan
física o geométricamente situaciones diferentes, sino la misma estructura
expresada con coordenadas distintas.
\subsection{Representaciones Irreducibles}
\noindent
Una representación \(R\) se dice reducible si deja invariante algún subespacio
no trivial del espacio vectorial sobre el que actúa. Es decir, si existe un
subespacio \(W\subset V\), con $W \neq \{0\}$, $W \neq V$, tal que
\begin{equation}
    D_R(g)W \subset W,
    \qquad
    \forall g\in G.
\end{equation}
\noindent
Esto significa que, si un vector empieza dentro de \(W\), la acción del grupo
lo mantiene siempre dentro de \(W\). El subespacio \(W\) no se mezcla con el
resto del espacio bajo la acción del grupo.
\medskip\newline
Si no existe ningún subespacio no trivial invariante, la representación se
llama irreducible.
\begin{equation}
    R \ \text{irreducible (irr)}
    \quad
    \Longleftrightarrow
    \quad
    \nexists\, W\subset V
    \ \text{no trivial tal que}\
    D_R(g)W \subset W.
\end{equation}
\noindent
Una representación es completamente reducible si puede escribirse como suma directa\footnote{Dado un espacio vectorial $V$, se dice que $V$ es la \textit{suma directa}
de dos subespacios $U$ y $W$, y se escribe $V = U \oplus W$, si todo vector
$v \in V$ se escribe de forma única como $v = u + w$ con $u \in U$ y $w \in W$,
o equivalentemente, si $V = U + W$ y $U \cap W = \{0\}$. Un elemento se escribe como un par ordenado $(u, w), \quad u \in U, \quad w \in W$, o equivalentemente como  $u + w.$
Ambas notaciones son equivalentes. La diferencia es que $u + w$ se emplea cuando
$U$ y $W$ son subespacios de un mismo espacio $V$, mientras que $(u, w)$ se utiliza
cuando $U$ y $W$ son espacios abstractos sin un ambiente común.} de representaciones irreducibles:
\begin{equation}
    R
    =
    R_1 \oplus R_2 \oplus \cdots \oplus R_n.
\end{equation}
\noindent
En una base adecuada, las matrices de la representación adoptan forma diagonal
por bloques\footnote{Cada bloque actúa sobre un subespacio invariante distinto.}:
\begin{equation}
    D_R(g)
    =
    \begin{pmatrix}
        D_{R_1}(g) & 0 & \cdots & 0 \\
        0 & D_{R_2}(g) & \cdots & 0 \\
        \vdots & \vdots & \ddots & \vdots \\
        0 & 0 & \cdots & D_{R_n}(g)
    \end{pmatrix}.
\end{equation}
\newpage
\noindent
\textbf{Ejemplo 1} 
\medskip\newline\noindent
Sea el grupo $G = \mathbb{Z}_2 = \{e, r\}$, con $r^2 = e$. Definimos la representación
$\rho : \mathbb{Z}_2 \to GL(\mathbb{R}^2)$ mediante:
\begin{equation}
    \rho(e) = \begin{pmatrix} 1 & 0 \\ 0 & 1 \end{pmatrix}, \qquad
    \rho(r) = \begin{pmatrix} 1 & 0 \\ 0 & -1 \end{pmatrix}.
\end{equation}
\noindent
El espacio $V = \mathbb{R}^2$ se descompone como $V = U \oplus W$, donde
\begin{align*}
    U &= \{(x, 0) : x \in \mathbb{R}\}, \\
    W &= \{(0, y) : y \in \mathbb{R}\}.
\end{align*}
\noindent
Ambos subespacios son invariantes bajo la acción del grupo:
\begin{equation}
    \rho(r)\begin{pmatrix} x \\ 0 \end{pmatrix} = \begin{pmatrix} x \\ 0 \end{pmatrix} \in U,
    \qquad
    \rho(r)\begin{pmatrix} 0 \\ y \end{pmatrix} = \begin{pmatrix} 0 \\ -y \end{pmatrix} \in W.
\end{equation}
\noindent
La restricción de $\rho$ a cada subespacio define dos representaciones irreducibles:
\begin{align*}
    \rho_1 &: \rho_1(e) = (1),\quad \rho_1(r) = (1)
    \qquad \text{(representación trivial)}, \\
    \rho_2 &: \rho_2(e) = (1),\quad \rho_2(r) = (-1)
    \qquad \text{(representación alternante)}.
\end{align*}
\noindent
Ambas son irreducibles por actuar sobre espacios de dimensión 1. La descomposición
$\rho = \rho_1 \oplus \rho_2$ se hace explícita en la estructura diagonal por bloques:
\begin{equation}
    \rho(r) = \begin{pmatrix} 1 & 0 \\ 0 & -1 \end{pmatrix}
    = \rho_1(r) \oplus \rho_2(r).
\end{equation}
\begin{flushright}
    \(\square\)
\end{flushright}
\noindent
\textbf{Ejemplo 2} 
\medskip\newline\noindent
Consideramos ahora la representación
$\rho : \mathbb{Z}_2 \to GL(\mathbb{R}^2)$ dada por:
\begin{equation}
    \rho(e) = \begin{pmatrix} 1 & 0 \\ 0 & 1 \end{pmatrix}, \qquad
    \rho(r) = \begin{pmatrix} 0 & 1 \\ 1 & 0 \end{pmatrix},
\end{equation}
\noindent
que corresponde a la matriz de permutación que intercambia las dos componentes.
Para encontrar los subespacios invariantes calculamos los autovalores de $\rho(r)$:
\begin{equation}
    \det(\rho(r) - \lambda I) = \lambda^2 - 1 = 0 \implies \lambda = \pm 1,
\end{equation}
\noindent
con autovectores asociados $v_1 = (1,1)$ para $\lambda = 1$ y $v_2 = (1,-1)$
para $\lambda = -1$. Los subespacios invariantes son entonces
\begin{align*}
    U &= \operatorname{span}\{(1,1)\}, \\
    W &= \operatorname{span}\{(1,-1)\},
\end{align*}
como se comprueba explícitamente:
\begin{equation}
    \rho(r)\begin{pmatrix}1\\1\end{pmatrix} = \begin{pmatrix}1\\1\end{pmatrix} \in U,
    \qquad
    \rho(r)\begin{pmatrix}1\\-1\end{pmatrix} = -\begin{pmatrix}1\\-1\end{pmatrix} \in W.
\end{equation}
\noindent
Tomando $\{v_1, v_2\}$ como nueva base, la matriz de cambio de base y su inversa son
\begin{equation}
    P = \begin{pmatrix} 1 & 1 \\ 1 & -1 \end{pmatrix}, \qquad
    P^{-1} = \frac{1}{2}\begin{pmatrix} 1 & 1 \\ 1 & -1 \end{pmatrix}.
\end{equation}
\noindent
La representación en la nueva base resulta
\begin{equation}
    \tilde{\rho}(r) = P^{-1}\rho(r)P =
    \frac{1}{2}\begin{pmatrix} 1 & 1 \\ 1 & -1 \end{pmatrix}
    \begin{pmatrix} 0 & 1 \\ 1 & 0 \end{pmatrix}
    \begin{pmatrix} 1 & 1 \\ 1 & -1 \end{pmatrix}
    = \begin{pmatrix} 1 & 0 \\ 0 & -1 \end{pmatrix},
\end{equation}
\noindent
que es diagonal por bloques. La descomposición en irreducibles es por tanto
\begin{equation}
    \tilde{\rho}(r) = \begin{pmatrix} 1 & 0 \\ 0 & -1 \end{pmatrix}
    = \rho_1(r) \oplus \rho_2(r),
\end{equation}
\noindent
donde $\rho_1$ es la representación trivial ($r \mapsto 1$) y $\rho_2$ es la
representación alternante ($r \mapsto -1$). La física no cambia al hacer este
procedimiento: únicamente hemos elegido la base en la que la estructura de suma
directa se hace manifiesta.
\section{Generadores de un grupo de Lie}
\noindent
En un grupo de Lie, cerca de la identidad, cualquier elemento puede escribirse en
términos de sus generadores. Si \(D_R\) es una representación del grupo, entonces, para una transformación
infinitesimal:
\begin{equation}
    D_R(\delta\theta)
    =
    \mathds{1}
    -
    i\delta\theta_a T^a_R
    +
    \mathcal{O}(\delta\theta^2).
\end{equation}
\noindent
donde los generadores \(T^a_R\) de la representación \(R\) se definen\footnote{El factor $i$ en la
definición se escoge de manera que, si en la representación $R$ los
generadores son hermíticos $\left(T^a_R\right)^\dagger = T^a_R.$, entonces las matrices $D_R(g)$ sean unitarias.
En este caso $R$ es una \textit{representación unitaria}.} como:
\begin{equation}
    T^a_R
    =
    i
    \frac{\partial D_R(\theta)}
    {\partial \theta_a}
    \bigg|_{\theta=0}.
\end{equation}
\noindent
En otras palabras, los generadores describen el comportamiento del grupo cerca
de la identidad. Para una transformación finita, la representación puede
escribirse como:
\begin{equation}
    D_R(\theta)
    =
    \exp
    \left(
        -i\theta_a T^a_R
    \right).
\end{equation}
\subsubsection{Propiedad 1}
\noindent
Toda representación $D_R$ unitaria es completamente reducible.
\subsubsection{Propiedad 2}
\noindent
Dadas dos matrices $
D_R(g_1)=\exp(-i\alpha_a T_R^a)$ y
$D_R(g_2)=\exp(-i\beta_a T_R^a),$ su producto es igual a $D_R(g_1g_2)$ y, por lo tanto, debe ser de la forma
$\exp(-i\delta_a T_R^a)$, para algún $\delta_a(\alpha,\beta)$,
\begin{equation}
e^{-i\alpha_a T_R^a} e^{-i\beta_a T_R^a}
=
e^{-i\delta_a T_R^a}.
\tag{2.6}
\end{equation}
\noindent
Si $A,B$ son matrices, en general
$e^A e^B \neq e^{A+B}$, por lo que en general
$\delta_a \neq \alpha_a+\beta_a$. Tomando el logaritmo y expandiendo hasta
segundo orden en $\alpha$ y $\beta$, obtenemos
\begin{equation}
\begin{aligned}
i\delta_a T_R^a
&=
\log\left\{
\left[
1+i\alpha_a T_R^a
+\frac{1}{2}(i\alpha_a T_R^a)^2
\right]
\left[
1+i\beta_a T_R^a
+\frac{1}{2}(i\beta_a T_R^a)^2
\right]
\right\}
\\
&=
\log\left[
1+i(\alpha_a+\beta_a)T_R^a
-\frac{1}{2}(\alpha_a T_R^a)^2
-\frac{1}{2}(\beta_a T_R^a)^2
-\alpha_a\beta_b T_R^a T_R^b
\right].
\end{aligned}
\tag{2.7}
\end{equation}
\noindent
Expandiendo el logaritmo, $\log(1+x)\simeq x-x^2/2$, y teniendo en cuenta
que los $T_R^a$ no conmutan, obtenemos:
\begin{equation}
\alpha_a \beta_b [T_R^a,T_R^b]
=
i\gamma_c(\alpha,\beta)T_R^c,
\tag{2.8}
\end{equation}
con:
\[
\gamma_c(\alpha,\beta)
=
-2\bigl(\delta_c(\alpha,\beta)-\alpha_c-\beta_c\bigr).
\]
Como esto debe cumplirse para todo $\alpha$ y $\beta$, $\gamma_c$ debe ser
lineal en $\alpha_a$ y en $\beta_a$, de modo que la relación entre $\gamma$
y $\alpha,\beta$ debe tener la forma general:
\[
\gamma_c=\alpha_a\beta_b f^{ab}{}_{c}
\]
para algunas constantes $f^{ab}{}_{c}$. 
\section{Álgebra de Lie del grupo}
\noindent
Cuando cambiamos la representación, en general la forma explícita e incluso las dimensiones de las matrices $D_R(g)$ cambiarán. Sin embargo, existe una propiedad importante de un grupo de Lie que es independiente de la
representación. Esta es su álgebra de Lie, que introducimos ahora.
\medskip\newline
Como hemos visto, los generadores de un grupo de Lie satisfacen relaciones de conmutación que
definen lo que llamamos álgebra de Lie del grupo. En general:
\begin{equation}
    \left[
        T^a_R,
        T^b_R
    \right]
    =
    i f^{abc} T^c_R.
\end{equation}
\noindent
Las constantes \(f^{abc}\) se llaman constantes de estructura del grupo. Son
independientes de la representación elegida y codifican la estructura local del
grupo cerca de la identidad.
\medskip\newline
Si todas las constantes de estructura son cero, $f^{abc}=0$, entonces todos los generadores conmutan, y el grupo es abeliano.
\medskip\newline
Así, las constantes de estructura definen el álgebra de Lie, y el problema
de encontrar todas las representaciones matriciales de un álgebra de Lie se
reduce al problema algebraico de encontrar todas las posibles soluciones
matriciales $T_R^a$ de la ec (2.29).
\subsection{Operadores de Casimir}
\noindent
En el estudio de las representaciones, los operadores de Casimir desempeñan
un papel muy importante. Estos son operadores construidos a partir de los $T^a$, por lo
que conmutan con todos los generadores del grupo. Es decir, un operador
\(C\) es de Casimir si
\begin{equation}
    \left[
        C,
        T^a
    \right]
    =
    0,
    \qquad
    \forall a.
\end{equation}
\noindent
Por el lema de Schur, en una representación irreducible todo operador que
conmuta con todos los generadores debe ser proporcional a la identidad:

\begin{equation}
    C = \lambda \mathds{1}.
\end{equation}
\noindent
De esta forma, podemos usar la constante \(\lambda\) para etiquetar representaciones irreducibles.
\medskip\newline
\textbf{Ejemplo 3} 
\medskip\newline\noindent
En \(SU(2)\), los generadores satisfacen
\begin{equation}
    \left[
        J_i,
        J_j
    \right]
    =
    i\epsilon_{ijk}J_k.
\end{equation}
El operador
\begin{equation}
    \vec{J}^{\,2}
    =
    J_1^2 + J_2^2 + J_3^2
\end{equation}
conmuta con todos los generadores:
\begin{equation}
    \left[
        \vec{J}^{\,2},
        J_i
    \right]
    =
    0.
\end{equation}
Por tanto, \(\vec{J}^{\,2}\) es un operador de Casimir. En una representación
irreducible de \(SU(2)\), se cumple
\begin{equation}
    \vec{J}^{\,2}
    =
    j(j+1)\mathds{1},
\end{equation}
donde \(j\) etiqueta la representación.
\subsection{Grupos compactos}
\noindent
Un grupo de Lie se dice que es compacto si su variedad paramétrica es compacta. De
forma intuitiva, esto significa que sus parámetros viven en un espacio cerrado
y acotado.
\medskip\newline
Por ejemplo, el grupo de rotaciones es compacto, ya que sus parámetros
angulares están acotados y las rotaciones vuelven periódicamente sobre sí
mismas. En cambio, el grupo de traslaciones espaciales no es compacto, porque
sus parámetros pueden tomar valores arbitrariamente grandes:
\medskip\newline
Si el grupo es compacto, entonces sus representaciones irreducibles de dimensión
finita son unitarias. Además, los parámetros que etiquetan las representaciones
suelen tomar valores discretos. En cambio, para grupos no compactos, pueden aparecer representaciones
unitarias de dimensión infinita y parámetros continuos. Por ello, las
representaciones de dimensión finita de grupos no compactos no son, en general,
unitarias.
\medskip\newline
La relevancia física se encuentra en el hecho de que en una
representación unitaria los generadores son operadores hermíticos y, de
acuerdo con las reglas de la mecánica cuántica, solo los operadores hermíticos
pueden identificarse con observables. Si un grupo es no compacto, para
identificar sus generadores con observables físicos necesitamos una
representación de dimensión infinita. Veremos en este capítulo que los grupos
de Lorentz y de Poincaré son no compactos, y que las representaciones de
dimensión infinita se obtienen introduciendo el espacio de Hilbert de estados
de una partícula.
\newpage
\section{Grupo de Lorentz}
\noindent

%------------------------------------------------------------------------------------------------------------------------------------------
\clearpage
\numberwithin{equation}{section}
\numberwithin{table}{section}
\numberwithin{figure}{section}

%\clearpage
%\section{Appendix - Raw data}

\end{document}

